Sleep is not a luxury but a biological necessity. Its chronic absence has been declared a public health problem by the CDC, with over one-third of U.S. adults failing to get the recommended 7-9 hours \citep{cdc2016shortsleep}. This deficit carries staggering socioeconomic costs:

\begin{itemize}
    \item \textbf{Economic Loss:} Insufficient sleep is estimated to cost the U.S. economy up to \$411 billion per year (about 2\% of GDP) in lost productivity from underperforming or absent workers \citep{rand2016cost}.
    \item \textbf{Public Safety:} Impaired attention and reaction time from fatigue are major dangers. In 2017, drowsy driving was estimated to have caused 91,000 police-reported crashes, 50,000 injuries, and 800 deaths in the U.S. alone \citep{nhtsa2017drowsy}.
    \item \textbf{Health Consequences:} Chronic short sleep is linked to a higher risk of hypertension, diabetes, obesity, heart disease, stroke, and depression. One study found that averaging under 6 hours of sleep per night increases the risk of all-cause mortality by about 10\% \citep{cdc2016shortsleep}.
\end{itemize}
